\chapter{}

“三洲田起义”的旧址如今是深圳东部华侨城的一处景点，一座小屋里几个蜡像，似乎在密谋着什么。他们在谈论什么呢？此时历史刚刚进入 20 世纪，这个古老帝国面临的问题比故宫城门上的铜钉还要多。只有一点是明确的，他们要组织一次暴动，推翻这个王朝。为首的人叫郑士良，也是刚刚被孙中山任命为司令。参加密谋的有当地绿林、三合会，以及广州洪门的首领。至于这次暴动能不能成功，将产生什么影响，他们都还来不及考虑。据孙中山发来的电报说，他们将向厦门发展，在那里会有来自台湾的军火支援。当革命史被用来消费的时候，当门票收入不能囊括这个景点的时候，大概只能告诉人们这么多。连同后来被水库淹没的“庚子革命首义中山纪念学校”一样，这些蜡像永远陷入了沉默。

其实，我更愿意相信，他们谈论的是失败。这个垂死的王朝太需要流血、牺牲、乃至失败来警醒世人了。其实，失败的阴影已经向他们迫近，谈论的话题也更加具体一些：这间小屋的主人要撵他们走。这里本来就是个卖杂货的“义合小铺”，做生意人气旺不怕，高谈阔论也可以，可你真的要把这里当做起义指挥部人家就不干了，你玩真的啊？廖氏家族的耆老们都不干了，要革命可以，在廖家的祠堂里革命不可以。另外，有情报说广东水师提督已经嗅到气味了，清军已经进驻沙湾一带。还有，孙中山第二次由香港登陆受阻，他已经随船开向了日本。孙中山的意思是：“如机密已泄，应暂行解散以避敌锋。”然而郑士良深知敌情，举事已箭在弦上，此时更不能遽离山寨，否则后果不堪设想。怎么办？干！

于是，真正的首义地点转移到了离三洲田十多公里的罗生大屋，这里也是另一位兴中会会员的祖屋。罗生大屋前的广场有半个足球场那么大，站了600 名参加烧猪祭旗的勇士们。他们有 300 条枪，每枪配子弹 30 发， 阵容也还可以。兴中会的唐梦尧先行揭礼：“众位兄弟，百打百胜，到来就位！” 众人高呼：“剑起灭匈奴，同申九世仇，汉人连处立，即日复神州！”“跟孙中山要跟到底！”等口号，宣告了起义开始。然后，互相剪辫子，缠红头巾红腰带，卷起裤脚，大口喝鸡血酒。这一脚高一脚低的裤脚日后就成为红头军的特有军容，而兴中会“驱除鞑虏，恢复中华，建立合众国”的口号也正式成为他们的行动纲领。

于是，1900 年10 月6 日，农历农历闰八月十三日，这个金风秋月乾坤朗朗的晚上，中国跨入了另一个时代。 为什么如此粗俗的仪式，如此狭隘的口号，会被后人重视？孙中山这样说：“经此次（指三洲田起义）失败而后，回顾中国之人心，已觉与前有别矣。当初次（指乙未广州起义）之失败也，举国舆论，莫不目予辈为乱臣贼子，大逆不道，咒诅谩骂之声，不绝于耳。吾人足迹所到，凡认识者，几视为毒蛇猛兽，而莫敢与吾人交游也，惟庚子失败之后，则鲜闻一般人之恶声相加，而有识之士，且多为吾人扼腕叹惜，恨其事之不成矣。前后相较，差若天渊。吾人睹此情形，中心快慰，不可言状。知国人之迷梦已有渐醒之兆。”

前后天渊之别只因有了明确的革命目标，有了蓝图，人家知道你要干什么要到哪里去。

在此前，广州起义的口号是“除暴安良”，在这同时及以后北方义和团运动的口号是“扶清灭洋”，都带着旧时代的印记，都让人觉得你们不过是想造反，造反成功你自己当皇帝，如果不成功你就是流寇盗贼，所以同情的人并不多。那时的孙中山一心想推翻满清政府。为达到目的他甚至跟日本人合作过，跟英国人合作过，还要跟李鸿章合作，对于建设一个什么样的国家并不清楚。经过那些失败，他的革命理想才逐渐清晰。这个时期大大小小的起义有很多次，也失败了很多次，只有这一次失败特别有价值。因为明确了要建立一个合众国，他要推翻的是一种制度，他要建设的是一种新型国家，国家的主人是老百姓自己。孙中山感慨道“有志之士多起救国之思，革命风潮自此萌芽。”

这其中有两个细节值得特别关注：剪辫子和卷裤脚。我们知道满清王朝对中国实现统治有一个三降三不降的承诺，即所谓“官降民不降，武降文不降，男降女不降”，而男人投降的标志就是剃头蓄发留辫子，所谓留发不留头，留头不留发。鲁迅小说《风波》即是对剪去辫子后国民心理的精彩描绘。所以杀猪祭旗的勇士们剪去辫子，则意味着对清王朝的彻底决裂。再一个是卷裤脚，这些勇士们不以整齐划一为美，反而故意把裤腿弄得一脚高一脚低，所为何来？细细品味，乃知正是造反者的美学诉求，庄稼汉泥腿子犯上作乱，蔑视旧秩序是也。往日衣冠楚楚克己复礼君君臣臣统统不算什么，老子才够威够猛！

这才是“现代性”的真义。现代性是指区别于古代国家的发展变化着的那些要素特征。它不仅是时间的，预示着历史的动向，引领着潮流浩浩荡荡；它更是空间的，意味着历史的多种可能性，人民有权选择不一样的活法。

建立合众国，让老百姓成为国家的主人，是让中国告别古代社会，跨入现代社会的第一声呐喊。庚子起义只延续了 32 天，但它打响了民主革命的发令枪，自此中国人民登上了历史舞台，演出了近百年威武雄壮的活剧，三洲田和罗氏大屋也成为中国跨入现代留下的第一只脚印。